% KYUSHU UNIVERSITY beamer template
% Created by Yu in 2024

\documentclass[11pt]{beamer}
% You can use \documentclass[11pt,aspectratio=169]{beamer} 
% to adjust the aspect ratio to 16:9.

\usepackage{graphicx}
\usepackage{listings}
\usepackage[
	backend=biber,
	style=verbose,
	sorting=ynt
]{biblatex}

\addbibresource{references.bib}

\usetheme{Madrid}

\author[Grupo 3]{Catalina Leal Rojas, Lucas Daniel Carrillo Aguirre, Lucas Eduardo Veras Costa, Maria Paula Basto Lozano}
\title[Sección 5]{Predicción del salario}

\setbeamertemplate{navigation symbols}{}


% Uncomment the section below to show a table of contents slide every time a new section starts
% \AtBeginSection[]
% {
%  \begin{frame}<beamer>
%  \frametitle{Table of contents}
%  \tableofcontents[currentsection]
%  \end{frame}
% }

\usepackage{Kyushu}

\begin{document}

\begin{frame}
\titlepage
\end{frame}

\section{Example}

\begin{frame}{Modelos seleccionados}
\begin{itemize}

\item \textbf{Modelo 1 (Edad cuadrática):}
{\small
\[
\begin{aligned}
\ln w_{i} &= \beta_{0} 
+ \beta_{1}\,\text{Edad}_{i}
+ \beta_{2}\,\text{Edad}_{i}^{2}
+ u_{i}
\end{aligned}
\]
}

\item \textbf{Modelo 2 (Brecha salarial por género incondicional):}
{\small
\[
\begin{aligned}
\ln w_{i} &= \beta_{0} 
+ \beta_{1}\,\text{Mujer}_{i}
+ u_{i}
\end{aligned}
\]
}

\item \textbf{Modelo 3 (Brecha salarial por género condicional):}
{\small
\[
\begin{aligned}
\ln w_{i} &= \beta_{0} 
+ \beta_{1}\,\text{Mujer}_{i}
+ \beta_{2}\,\text{Edad}_{i}
+ \beta_{3}\,\text{Edad}_{i}^{2}
+ \beta_{4}\,\text{Educ}_{i}
\\&\quad
+ \beta_{5}\,\text{Horas}_{i}
+ \beta_{6}\,\text{Microemp}_{i}
+ \beta_{7}\,\text{Formal}_{i}
+ \beta_{8}\,\text{TamañoFirma}_{i}
\\&\quad
+ \beta_{9}\,\text{Oficio}_{i}
+ u_{i}
\end{aligned}
\]
}
\end{itemize}


\end{frame}

\begin{frame}{Modelos seleccionados}

\begin{itemize}
\item \textbf{Modelo 4 (Modelo 3 con polinomio de tercer grado en edad e interacciones de este polinomio con el resto de controles):}
{\tiny
\[
\begin{aligned}
\ln w_{i} &= \beta_{0} 
+ \sum_{m=1}^{3} \beta_{1}^{m}\,\text{Edad}_{i}^{m}
+ \beta_{2}\,\text{Mujer}_{i}
+ \beta_{3}\,\text{Educ}_{i}
+ \beta_{4}\,\text{Horas}_{i}
\\&\quad
+ \beta_{5}\,\text{Microemp}_{i}
+ \beta_{6}\,\text{Formal}_{i}
+ \beta_{7}\,\text{TamañoFirma}_{i}
+ \beta_{8}\,\text{Oficio}_{i}
\\&\quad
+ \sum_{m=0}^{3}\Big[
\gamma_{1}^{m}(\text{Mujer}_{i}\,\text{Edad}_{i}^{m})
+ \gamma_{2}^{m}(\text{Educ}_{i}\,\text{Edad}_{i}^{m})
+ \gamma_{3}^{m}(\text{Horas}_{i}\,\text{Edad}_{i}^{m})
\\&\qquad\qquad
+ \gamma_{4}^{m}(\text{Microemp}_{i}\,\text{Edad}_{i}^{m})
+ \gamma_{5}^{m}(\text{Formal}_{i}\,\text{Edad}_{i}^{m})
+ \gamma_{6}^{m}(\text{TamañoFirma}_{i}\,\text{Edad}_{i}^{m})
\\&\qquad\qquad
+ \gamma_{7}^{m}(\text{Oficio}_{i}\,\text{Edad}_{i}^{m})
\Big]
+ u_{i}
\end{aligned}
\]
}

\item \textbf{Modelo 5 (Modelo 4 pero con el polinomio de edad de grado 5):}
{\small
\[
\begin{aligned}
\ln w_{i} &= \beta_{0} 
+ \sum_{m=1}^{5} \beta_{1}^{m}\,\text{Edad}_{i}^{m}
+ \sum_{j} \beta_{j}\,\text{Control}_{ji}
\\&\quad
+ \sum_{m=0}^{5}\sum_{j} \gamma_{j}^{m}\,(\text{Control}_{ji}\,\text{Edad}_{i}^{m})
+ u_{i}
\end{aligned}
\]
}
\emph{Controles:} Mujer, Educ, Horas, Microemp, Formal, TamañoFirma, Oficio.

\end{itemize}
    
\end{frame}

\begin{frame}{Modelos seleccionados}
\begin{itemize}

\item \textbf{Modelo 6 (Modelo 3 añadiendo Eduación Terciaria como control):}
{\small
\[
\begin{aligned}
\ln w_{i} &= \beta_{0} 
+ \beta_{1}\,\text{Mujer}_{i}
+ \beta_{2}\,\text{Edad}_{i}
+ \beta_{3}\,\text{Edad}_{i}^{2}
+ \beta_{4}\,\text{Educ}_{i}
\\&\quad
+ \beta_{5}\,\text{Horas}_{i}
+ \beta_{6}\,\text{Microemp}_{i}
+ \beta_{7}\,\text{Formal}_{i}
+ \beta_{8}\,\text{Tamaño}_{i}
+ \beta_{9}\,\text{Oficio}_{i}
\\&\quad
+ \beta_{10}\,\text{College}_{i}
+ u_{i}
\end{aligned}
\]
}

\item \textbf{Modelo 7 (Modelo 6 con polinomios de grado 5 para edad y total de horas trabajadas):}
{\small
\[
\begin{aligned}
\ln w_{i} &= \beta_{0} 
+ \sum_{m=1}^{5} \beta_{1}^{m}\,\text{Edad}_{i}^{m}
+ \sum_{m=1}^{5} \beta_{2}^{m}\,\text{Horas}_{i}^{m}
\\&\quad
+ \beta_{3}\,\text{Mujer}_{i}
+ \beta_{4}\,\text{Educ}_{i}
+ \beta_{5}\,\text{Microemp}_{i}
+ \beta_{6}\,\text{Formal}_{i}
+ \beta_{7}\,\text{Tamaño}_{i}
\\&\quad
+ \beta_{8}\,\text{Oficio}_{i}
+ \beta_{9}\,\text{College}_{i}
+ u_{i}
\end{aligned}
\]
}
  \end{itemize}  
\end{frame}

\begin{frame}{Modelos seleccionados}
\begin{itemize}

\item \textbf{Modelo 8 (Modelo 6 con polinomios de grado 3 para edad y educación, así como interacciones de estos polinomios con el resto de controles):}
{\small
\[
\begin{aligned}
\ln w_{i} &= \beta_{0} 
+ \sum_{m=1}^{3} \beta_{1}^{m}\,\text{Edad}_{i}^{m}
+ \sum_{m=1}^{3} \beta_{2}^{m}\,\text{Horas}_{i}^{m}
\\&\quad
+ \sum_{m=0}^{3}\sum_{j} \gamma_{1j}^{m}\,(\text{Edad}_{i}^{m}\,\text{Control}_{ji})
+ \sum_{m=0}^{3}\sum_{j} \gamma_{2j}^{m}\,(\text{Horas}_{i}^{m}\,\text{Control}_{ji})
+ u_{i}
\end{aligned}
\]
}
\emph{Controles:} Mujer, Educ, Microemp, Formal, Tamaño, Oficio, College.

\end{itemize}
    
\end{frame}

\begin{frame}{Resultados RMSE}
\begin{table}[htbp]
\centering
\begin{tabular}{l r}
\hline
Modelo & RMSE \\
\hline
1 & 0.67766 \\
2 & 0.69275 \\
3 & 0.44956 \\
4 & 0.44769 \\
5 & 0.54028 \\
6 & 0.44956 \\
7 & 0.44763  \\
8 & 0.62034 \\
\hline
\end{tabular}
\caption{Comparación del RMSE entre los modelos seleccionados}
\label{tab:scores}
\end{table}
    
\end{frame}

\begin{frame}{Resultados RMSE}

\begin{itemize}
    \item A medida que incluimos más controles y formas funcionales flexibles, el
RMSE disminuye de manera importante, lo que evidencia la relevancia de capturar la heterogeneidad
de las variables individuales y laborales para disminuir el sesgo.

\item Los modelos 4 (polinomio de grado 3 en edad e interacciones con todos los controles) y 7 (polinomio de grado 5 en edad y horas trabajadas) tuvieron el menor RMSE, aunque estos valores no están lejos de los RMSE correspondientes a los modelos 3 y 6.

\item Esto indica que gran parte de la capacidad predictiva ya está
capturada por los controles, pero que la incorporación de estas formas funcionales adicionales ayuda a
refinar la predicción en los márgenes donde las relaciones son más complejas.

\item Los modelos 5 y 8 son un ejemplo de overfitting.
\end{itemize}
   
\end{frame}

\begin{frame}{Análisis de los errores en el modelo 7}


\begin{figure}[htbp]
    \centering
    \includegraphics[width=0.8\textwidth]{hist_error_p5.png}
    \caption{Histograma del error de predicción en el "testing set"}
    \label{fig:hist_error}
\end{figure}
    
\end{frame}

\begin{frame}{Análisis de los errores en el modelo 7}

\begin{itemize}
    \item Distribución centrada en 0 y simétrica: el modelo no presenta un sesgo sistemático, este predice tanto por encima como por debajo del valor observado con una frecuencia similar.
    \item Para las observaciones con errores más altos (percentiles 1 y 99 de la distribución) tienden a registrar menos horas trabajadas ( 43.97 en promedio, lo cual está por debajo del percentil 25 en la muestra total), menor grado de formalidad y una mayor presencia de trabajadores en microempresas. Esto sugiere que los errores más grandes no provienen de perfiles totalmente atípicos, sino de segmentos del mercado laboral donde la variabilidad en los ingresos es mayor y más difícil de capturar con el modelo. 
\end{itemize}
    
\end{frame}

\begin{frame}{LOOCV (modelos 4 y 7)}
\begin{itemize}
    
\item Metodología válida para regresión polinomial:
\end{itemize}
\[
CV_{(n)} = \frac{1}{n} \sum_{i=1}^{n} \left( \frac{y_i - \hat{y}_i}{1 - h_i} \right)^{2},
\]
    
\end{frame}

\begin{frame}{LOOCV (modelos 4 y 7)}

\begin{table}[H]
\centering
\begin{tabular}{lcc}
\hline
                 & Modelo 4 & Modelo 7 \\
\hline
RMSE VS          &    0.4477      &    0.4476       \\
RMSE LOOCV       &    0.4440      &    0.4504      \\
\hline
\end{tabular}
\caption{Comparación de errores de predicción}
\end{table}
    
\end{frame}

\begin{frame}{LOOCV (modelos 4 y 7)}
\small
\begin{itemize}
    \item El hecho de que el Modelo 4 mantenga un error estable e incluso algo menor bajo LOOCV sugiere que es más robusto y generaliza mejor a nuevas observaciones. En contraste, el Modelo 7 parece más sensible al esquema de validación.
    \item En el modelo 4 la inclusión de un menor grado en el polinomio (3 vs 5) y además interacciones con los controles permite incluir heterogeneidad de forma más estructurada.
    \item Por su parte, el Modelo 7 introduce polinomios de quinto grado, lo que genera curvas más sensibles y con oscilaciones más marcadas, en especial en los extremos del soporte. Al no estar acompañada de interacciones, esta flexibilidad resulta menos organizada y más vulnerable a la exclusión de observaciones. 
    \item Es probable que el Modelo 7 resulta más sensible a observaciones influyentes, lo que eleva su RMSE con LOOCV, mientras que el Modelo 4 muestra mayor estabilidad frente a outliers o alto apalancamiento.
\end{itemize}
    
\end{frame}


\end{document}
